\FloatBarrier
\section{چرا \lr{v2}}
\label{sec:why}

\lr{v2} برای رفع کاستی‌های بخش~\ref{sec:gaps} طراحی شد. مهم‌ترین تغییر این است که مدل به‌جای یک خروجی، سه خروجی جدا دارد:
\begin{itemize}
\item \textbf{\lr{K}} فقط از سینماتیک و تاریخچه‌ی بلندکردن‌ها ساخته می‌شود. می‌گوید سینماتیک به‌تنهایی چقدر خوب پیش‌بینی می‌کند.
\item \textbf{\lr{E}} فقط از \lr{EEG} ساخته می‌شود. می‌گوید \lr{EEG} به‌تنهایی چه چیزی از فعالیت عضله را در خود دارد.
\item \textbf{\lr{F}} همان \lr{K} است به‌علاوه‌ی یک اصلاح کوچک از \lr{EEG}. اگر \lr{F} از \lr{K} بهتر شود، \lr{EEG} چیزی دارد که سینماتیک ندارد.
\end{itemize}
با این سه خروجی، دو پرسشی که در مدل پایه با هم قاطی می‌شدند جدا می‌شوند. جدول~\ref{tab:why} بقیه‌ی تغییرها را کنار کاستی متناظر می‌گذارد.

\begin{table}[htbp]
\centering\small
\caption{کاستی‌های مدل پایه و پاسخ \lr{v2} به هر کدام.}
\label{tab:why}
\begin{tabular}{p{0.38\textwidth}p{0.52\textwidth}}
\toprule
کاستی مدل پایه & تصمیم در \lr{v2} \\
\midrule
یک خروجی؛ سهم \lr{EEG} را نمی‌توان جدا کرد & سه خروجی جدا: \lr{K} (سینماتیک)، \lr{E} (فقط \lr{EEG})، \lr{F} (\lr{K} به‌علاوه‌ی اصلاح \lr{gated} از \lr{EEG})؛ \lr{K} از \lr{E} و \lr{F} هیچ گرادیانی نمی‌گیرد \\
پنجره‌ی یک‌ثانیه‌ای & \lr{tokens} ۵۰ هرتز و پنجره‌ی ۸ ثانیه‌ای (۴۰۰ \lr{tokens}) \\
بی‌خبری از بلندکردن‌های قبلی & هشت ویژگی تاریخچه از نیروها، فقط از گذشته و بدون \lr{EMG} \\
\lr{loss function} هشت‌جزئی و یک‌طرفه در حوزه‌ی $z$ & حوزه‌ی جذر دامنه و \lr{loss function} چهارجزئی، بدون وزن اضافه‌ی قله و جزء استراحت \\
\lr{EEG} خام در حوزه‌ی زمان & \lr{band power} در چهار باند برای هر الکترود، به‌صورت گرافی از الکترودها \\
اطلاعات از چند نوع لایه جابه‌جا می‌شود؛ \lr{rollout} تعریف‌پذیر نیست & قاعده‌ی «یک مخلوط‌کننده»: فقط لایه‌های توجه اطلاعات را بین موقعیت‌ها جابه‌جا می‌کنند و همه‌شان ثبت می‌شوند \\
همه‌چیز با هم آموزش می‌بیند و با هم متوقف می‌شود & هر جریان \lr{optimizer} و توقف خودش را دارد؛ جریانی که دیگر بهتر نمی‌شود در بهترین نقطه‌اش \lr{freeze} می‌شود \\
\bottomrule
\end{tabular}
\end{table}

\lr{v2} از سه جهت بهتر از مدل پایه است.
\begin{enumerate}
\item \textbf{دقت.} \lr{K} در هر ۱۲ شرکت‌کننده \lr{VAF} بالاتری از مدل پایه دارد (بخش~\ref{sec:v2-results}).
\item \textbf{پرسش علمی روشن.} «آیا \lr{EEG} چیزی بیش از سینماتیک دارد؟» با مقایسه‌ی \lr{F} و \lr{K} و با \lr{residual test} جواب می‌گیرد. «آیا \lr{EEG} به‌تنهایی حرکت را رمز می‌کند؟» با \lr{E} جواب می‌گیرد.
\item \textbf{تفسیرپذیری.} برای هر عدد خروجی می‌توان دنبال کرد که از کدام ورودی‌ها و از کدام لایه‌ها آمده است.
\end{enumerate}

\FloatBarrier
\section{مدل \lr{v2}}
\label{sec:v2}

شکل~\ref{fig:v2-arch} کل مدل را نشان می‌دهد. مدل دو «جریان» مستقل دارد که هر کدام یک خروجی می‌سازند: جریان سینماتیک \lr{K} را و جریان \lr{EEG} \lr{E} را. یک ماژول کوچک سوم (\lr{fusion}) این دو را برای ساختن \lr{F} به هم وصل می‌کند. در ادامه اول ورودی‌ها و خروجی‌ها، بعد اجزای مدل و در آخر نحوه‌ی آموزش توضیح داده می‌شود.

\subsection{ورودی‌ها، \lr{tokens} و خروجی‌ها}

همه‌ی ورودی‌ها از پیش‌پردازش بخش~\ref{sec:preproc-v2} می‌آیند؛ اینجا فقط خلاصه‌شان آمده است.

\textbf{نرخ \lr{token}.} مدل پایه با ۵۰۰ نمونه در ثانیه کار می‌کرد. برای اینکه \lr{v2} بتواند ۸ ثانیه را یک‌جا ببیند، همه‌ی سیگنال‌ها با فیلتر ضدتداخل به ۵۰ هرتز برده می‌شوند. هر \lr{token} نماینده‌ی ۲۰ میلی‌ثانیه است و یک پنجره‌ی ۸ ثانیه‌ای ۴۰۰ \lr{token} دارد.

\textbf{سینماتیک} $X^{k}\in\mathbb{R}^{T\times24}$. همان ۲۴ ویژگی مدل پایه (وضعیت دست، دو نیروی \lr{gated} و سرعت‌هایشان)، استانداردشده با آماره‌های آموزش.

\textbf{تاریخچه} $X^{h}\in\mathbb{R}^{T\times8}$. این ویژگی‌ها به مدل می‌گویند تا این لحظه در سری چه گذشته است: آیا و چند بار جسم نگه داشته شده، چقدر از شروع سری و از آخرین نگه‌داشتن گذشته، و نیروهای آخرین نگه‌داشتن و میانگین همه‌ی نگه‌داشتن‌های قبلی. مدل فقط وقتی از یک نگه‌داشتن خبردار می‌شود که آن نگه‌داشتن تمام شده باشد؛ پس هیچ ویژگی‌ای از آینده خبر ندارد.

\textbf{\lr{EEG}} $X^{e}\in\mathbb{R}^{T\times32\times4}$. به‌جای موج خام، برای هر الکترود لگاریتم توان در چهار باند: تتا (۴ تا ۸ هرتز)، آلفا (۸ تا ۱۳)، بتا (۱۳ تا ۳۰) و گامای پایین (۳۰ تا ۴۰)، استانداردشده و بریده در $\pm5$.

\textbf{هدف در حوزه‌ی جذر.} برای هر عضله $s_c$ صدک ۹۵ پوش در داده‌ی آموزش است. مدل به‌جای خود پوش، جذر پوش نرمال‌شده را پیش‌بینی می‌کند:
\begin{equation}
u=\sqrt{y/s_c},\qquad \hat y=s_c\,\hat u^{2}.
\end{equation}
جذر قله‌های بلند را کوتاه‌تر و سطح‌های پایین را بازتر می‌کند. در نتیجه خطای قله‌ها و خطای استراحت وزن نزدیک‌تری می‌گیرند و دیگر لازم نیست قله‌ها را جداگانه وزن بدهیم یا جزء جداگانه‌ای برای استراحت بگذاریم.

\textbf{خروجی‌ها.} مدل برای هر \lr{token} و هر عضله سه عدد برمی‌گرداند:
\begin{equation}
K=K_{base}\odot g(X^{h}),\qquad E=\mathrm{E\text{-}stream}(X^{e}),\qquad F=\mathrm{relu}\big(K+F_{eeg}\big).
\end{equation}
$K_{base}$ پیش‌بینی جریان سینماتیک است و $g$ یک ضریب از تاریخچه که آن را کمی بزرگ یا کوچک می‌کند. $F_{eeg}$ اصلاحی است که \lr{EEG} به \lr{K} اضافه می‌کند. مدل این اجزا را هم جداگانه برمی‌گرداند تا بتوان سهم هر کدام را دید.

\begin{figure}[p]
\centering
% KG-GT v2: overview in the style of a paper's first figure (portrait).
\begin{latin}\resizebox{\textwidth}{!}{%
\begin{tikzpicture}[font=\footnotesize, box/.append style={font=\scriptsize}]
% ================= inputs
\node[lab, font=\normalsize\bfseries, anchor=west] at (-0.3,0.65) {Inputs: one window of 8 s = 400 tokens at 50 Hz (anti-aliased reduction from 500 Hz)};
\node[lab, font=\small\bfseries, anchor=west] at (-0.3,0.0) {Kinematics};
\node[lab, anchor=west] at (1.95,0.0) {24 features (as in the baseline)};
\foreach \k in {0,...,3}{\begin{scope}[shift={(0,-0.45-0.3*\k)}]\sig{0}{5.3}{0.1}{kinc!85!black}{70*\k+15}\end{scope}}
\node[lab, font=\small\bfseries, anchor=west] at (-0.3,-1.85) {Trial history};
\node[lab, anchor=west] at (2.25,-1.85) {8 features, past only, no EMG};
\draw[kinc!60!black, line width=0.8pt] (0,-2.8)--(1.1,-2.8)--(1.1,-2.55)--(2.6,-2.55)--(2.6,-2.4)--(4.2,-2.4)--(4.2,-2.3)--(5.3,-2.3);
\node[lab, anchor=west, align=left] at (-0.3,-3.25) {holds so far, time since last hold, previous and mean load / grip force};
\node[lab, font=\small\bfseries, anchor=west] at (9.6,0.0) {EEG band power};
\headgraph{h0}{10.75,-1.45}{1.05}
\node[lab, anchor=west, align=left] at (12.05,-1.45) {32 electrodes $\times$ 4 bands\\ ($\theta$ 4--8, $\alpha$ 8--13, $\beta$ 13--30,\\ low $\gamma$ 30--40 Hz)\\ log power, z-scored, clip $\pm$5};

% ================= kinematic stream (x = 2.4)
\node[kbox, align=center, text width=4.3cm] (stem) at (2.4,-4.35) {conv stem, kernel 5 (100 ms)\\ 32 $\to$ 64 channels + sinusoidal positions};
\draw[arr] (2.4,-3.5)--(stem.north);
\tblock{kenc}{2.4,-7.0}{4.3cm}{Kinematic temporal encoder}{3}{kinc}{LayerNorm $\to$ attention $\to$ add}{LayerNorm $\to$ FFN $\to$ add}{pre-norm block, d = 64, 4 heads, FFN 256}
\draw[arr] (stem.south)--(kenc.north);
\node[kbox, align=center, text width=4.3cm] (kproj) at (2.4,-9.15) {each token $\to$ 5 muscle nodes (32-d)};
\draw[arr] (kenc.south)--(kproj.north);
\musclegraph{km}{2.4,-11.2}{1.1}{kinc}
\draw[arr] (kproj.south)--(kmAD.north);
\node[lab, anchor=east, align=right] at (0.8,-10.55) {muscle\\ graph $\times$2\\ prior: EMG\\ co-activation};
\node[kbox, align=center, text width=4.3cm] (kdec) at (2.4,-13.2) {per-muscle decoder + softplus $\to K_{base}$};
\draw[arr] (2.4,-12.35)--(kdec.north);
\node[circle, draw, inner sep=1pt, fill=white, font=\normalsize] (times) at (2.4,-14.3) {$\times$};
\draw[arr] (kdec.south)--(times.north);
\node[kbox, align=center, text width=1.9cm] (gain) at (5.4,-14.3) {history gain\\ MLP 8$\to$16$\to$5\\ $g\in[0.5,2]$};
\draw[arr, kinc!70!black] (5.4,-3.5)--(gain.north);
\draw[arr] (gain.west)--(times.east);

% ================= EEG stream (x = 13.0)
\node[ebox, align=center, text width=4.3cm] (emb) at (13.0,-4.35) {electrode embedding: MLP 4$\to$32\\ + identity vector per electrode};
\draw[arr] (10.75,-2.65)--(10.75,-3.25)-|(emb.north);
\headgraph{h1}{13.0,-6.65}{1.05}
\draw[arr] (emb.south)--(13.0,-5.38);
\node[lab, anchor=west, align=left] at (14.2,-6.65) {electrode\\ graph $\times$2,\\ per token\\ prior: scalp\\ distance};
\node[trapezium, trapezium left angle=70, trapezium right angle=70, shape border rotate=180, draw=eegc!70!black, fill=eegc!10, align=center, minimum width=3.8cm, font=\scriptsize] (pool) at (13.0,-8.4) {attention pooling over 32 electrodes};
\draw[arr] (13.0,-7.72)--(pool.north);
\tblock{eenc}{13.0,-10.6}{4.3cm}{EEG temporal encoder}{2}{eegc}{LayerNorm $\to$ attention $\to$ add}{LayerNorm $\to$ FFN $\to$ add}{pre-norm block, d = 64, 4 heads, FFN 256}
\draw[arr] (pool.south)--(eenc.north);
\node[ebox, align=center, text width=4.3cm] (eproj) at (13.0,-12.75) {each token $\to$ 5 muscle nodes (32-d)};
\draw[arr] (eenc.south)--(eproj.north);
\musclegraph{em}{13.0,-14.8}{1.1}{eegc}
\draw[arr] (eproj.south)--(emAD.north);
\node[lab, anchor=west, align=left] at (14.55,-13.75) {muscle\\ graph $\times$2};
\node[ebox, align=center, text width=4.3cm] (edec) at (13.0,-16.65) {per-muscle decoder + softplus};
\draw[arr] (13.0,-15.95)--(edec.north);

% ================= fusion (x = 8.2)
\node[fbox, align=center, text width=3.1cm] (cross) at (8.2,-10.6) {\textbf{cross-attention}\\ query: kinematic tokens\\ key, value: EEG tokens};
\draw[darr, kinc!70!black] ($(kenc.east)+(0.55,-0.8)$)-|node[lab, pos=0.25, above]{stop-gradient}(cross.north);
\draw[arr, eegc!70!black] (eenc.west)--(cross.east);
\node[fbox, align=center, text width=3.1cm] (gate) at (8.2,-12.7) {linear per muscle $\times$ gate $\gamma$\\ (starts at 0) $\to F_{eeg}$};
\draw[arr] (cross)--(gate);
\node[circle, draw, inner sep=1pt, fill=white, font=\normalsize] (plus) at (8.2,-15.3) {$+$};
\draw[arr] (gate)--(plus);
\coordinate (kt) at (2.4,-15.3);
\draw[darr, kinc!70!black] (kt)--node[lab, pos=0.62, below]{K, stop-gradient}(plus.west);
\node[fbox, align=center] (relu) at (8.2,-16.35) {$F=\mathrm{relu}(K+F_{eeg})$};
\draw[arr] (plus)--(relu);

% ================= outputs
\node[result, draw=kinc!80!black, fill=kinc!12, align=center, text width=3.2cm] (K) at (2.4,-17.75) {\textbf{K}: kinematics-led};
\node[result, draw=fusc!80!black, fill=fusc!12, align=center, text width=3.2cm] (F) at (8.2,-17.75) {\textbf{F}: fused};
\node[result, draw=eegc!80!black, fill=eegc!12, align=center, text width=3.2cm] (E) at (13.0,-17.75) {\textbf{E}: EEG only};
\draw[arr] (times.south)--(K.north); \draw[arr] (relu)--(F); \draw[arr] (edec)--(E);
\node[lab, align=center] at (8.2,-18.45) {EMG envelope of 5 muscles in the model domain $u=\sqrt{y/s_c}$, interpolated back to 500 Hz for scoring};
\begin{scope}[shift={(0.65,-19.6)}]\env{0}{3.5}{0.75}{kinc!85!black}{10}\end{scope}
\begin{scope}[shift={(6.45,-19.6)}]\env{0}{3.5}{0.75}{fusc!85!black}{10}\end{scope}
\begin{scope}[shift={(11.25,-19.6)}]\env{0}{3.5}{0.6}{eegc!85!black}{-10}\end{scope}

% ================= loss, training, recorder
\node[box, font=\footnotesize, fill=black!3, anchor=north west, align=left, text width=16.4cm] (loss) at (-0.3,-19.9) {\textbf{Loss} $L=L_K+L_E+L_F$, for each output: Huber ($\delta=0.5$) + 0.3$\,(1-\mathrm{CCC})$ + 0.5 level error per detected hold + 0.1 temporal difference.\\
\textbf{Training}: one AdamW for the kinematic stream, one for the EEG stream and fusion (gate without weight decay); a stream with no improvement for 30 epochs returns to its best weights and is frozen, the other one goes on.\\
\textbf{Attention recorder} (outside the model): all 13 attention layers, per head and with gradients $\to$ rollout, attention flow, gradient-weighted rollout.};
\end{tikzpicture}}\end{latin}

\caption{نمای کلی \lr{v2} از ورودی تا سه خروجی. چپ: جریان سینماتیک (\lr{stem} کانولوشنی، سه بلوک \lr{Transformer} زمانی، گراف پنج عضله، \lr{decoder} و \lr{history gain}) که \lr{K} را می‌سازد. راست: جریان \lr{EEG} (\lr{embedding} الکترودها، گراف ۳۲ الکترود روی سر، \lr{attention pooling}، دو بلوک \lr{Transformer} زمانی، \lr{muscle graph}) که \lr{E} را می‌سازد. وسط: \lr{fusion} با \lr{cross-attention} و یک \lr{gate} که از صفر شروع می‌شود و \lr{F} را می‌سازد. خط‌چین‌ها جاهایی‌اند که گرادیان قطع می‌شود: \lr{fusion} \lr{K} و \lr{tokens} سینماتیک را می‌خواند اما به آن‌ها گرادیان نمی‌فرستد، پس \lr{K} فقط با خطای خودش آموزش می‌بیند. کادر پایین \lr{loss function}، آموزش و \lr{attention recorder} را خلاصه می‌کند.}
\label{fig:v2-arch}
\end{figure}

\subsection{بلوک توجه یکسان}

همه‌ی لایه‌های توجه \lr{v2}، چه روی زمان، چه روی الکترودها و چه روی عضله‌ها، یک شکل دارند. هر موقعیت (یک \lr{token}، یک الکترود یا یک عضله) به بقیه‌ی موقعیت‌ها «نگاه می‌کند» و از هر کدام به اندازه‌ی وزن توجهش اطلاعات می‌گیرد. ماتریس توجه $A$ همین وزن‌ها را نگه می‌دارد: سطر $i$ موقعیتی است که اطلاعات می‌گیرد، ستون $j$ موقعیتی است که اطلاعات می‌دهد و جمع هر سطر ۱ است. برای سر $h$:
\begin{equation}
A^{h}=\operatorname{softmax}\!\Big(\frac{Q^{h}K^{h\top}}{\sqrt{d_h}}+M^{h}\Big),\qquad
x\leftarrow x+W_o\big[A^{1}V^{1}\,\|\,\dots\,\|\,A^{H}V^{H}\big],\qquad x\leftarrow x+\mathrm{FFN}(\mathrm{LN}(x)).
\end{equation}
$Q$، $K$ و $V$ از $\mathrm{LN}(x)$ ساخته می‌شوند (\lr{pre-norm}). «$x\leftarrow x+\dots$» یعنی خروجی هر بخش به ورودی‌اش اضافه می‌شود (\lr{residual connection})؛ پس اطلاعات هر موقعیت همیشه یک راه مستقیم هم به لایه‌ی بعد دارد.

$M$ یک \lr{bias} است که پیش از \lr{softmax} اضافه می‌شود. در توجه زمانی $M=0$ است. در لایه‌های گرافی $M^{h}=B^{h}+\lambda P$ است: $B^{h}$ یک \lr{bias} آموختنی برای هر یال در هر سر، $P$ یک ماتریس ثابت که دانش قبلی را وارد می‌کند (مثلاً اینکه الکترودهای نزدیک به هم مرتبط‌ترند) و $\lambda$ ضریبی که مدل یاد می‌گیرد چقدر به این دانش تکیه کند. دو لایه \lr{residual connection} ندارند: \lr{attention pooling} الکترودها و \lr{cross-attention} در \lr{fusion}.

\paragraph{قاعده‌ی «یک مخلوط‌کننده».} این قاعده \lr{v2} را تفسیرپذیر می‌کند. اطلاعات فقط از دو راه می‌تواند از یک موقعیت (زمان، الکترود، عضله یا جریان) به موقعیت دیگر برود: ماژول توجه بالا و ماژول \lr{attention pooling}. همه‌ی لایه‌های دیگر هر موقعیت را جدا پردازش می‌کنند. پس اگر وزن‌های توجه همه‌ی لایه‌ها را داشته باشیم، کل مسیر اطلاعات را داریم.

تنها استثنا کانولوشن ۵ضربه‌ای \lr{stem} سینماتیک است. ردپای آن به‌صورت یک ماتریس $S$ ثبت می‌شود: هر سطر جمع ۱ دارد و وزن هر جابه‌جایی برابر انرژی هسته‌ی کانولوشن در آن جابه‌جایی است.

وزن‌های توجه را یک ثبت‌کننده‌ی بیرونی (\lr{recorder}) برمی‌دارد: \lr{with recording(model) as tape}. این ثبت‌کننده توجه هر سر را پس از \lr{softmax} و پیش از \lr{dropout} نگه می‌دارد و در صورت نیاز گرادیان آن را هم نگه می‌دارد. ثبت کردن خروجی مدل را تغییر نمی‌دهد و لازم نیست به تابع \lr{forward} هیچ آرگومان اضافه‌ای داده شود. فهرست ۱۳ لایه‌ی ثبت‌شده از \lr{model.attention\_layers()} می‌آید.

\subsection{جریان سینماتیک}

این جریان از سینماتیک و تاریخچه، \lr{K} را می‌سازد:
\begin{enumerate}
\item \textbf{\lr{stem}.} یک کانولوشن یک‌بعدی با هسته‌ی ۵ روی سینماتیک و تاریخچه با هم (۳۲ ویژگی) که ۶۴ کانال می‌سازد. \lr{sinusoidal positional encoding} جای هر \lr{token} در پنجره را اضافه می‌کند.
\item \textbf{\lr{temporal encoder}.} ۳ بلوک توجه روی هر ۴۰۰ \lr{token} پنجره، با ۴ سر، \lr{feed-forward} ۲۵۶ و \lr{dropout} \lr{۰٫۲}. اینجا هر \lr{token} می‌تواند به هر لحظه‌ی دیگر پنجره نگاه کند.
\item \textbf{سر عضله.} هر \lr{token} به ۵ گره (یکی برای هر عضله، با عرض ۳۲) تبدیل می‌شود. دو بلوک توجه بین گره‌ها پیام جابه‌جا می‌کنند؛ ماتریس ثابت $P$ این گراف از هم‌فعالی عضله‌ها در \lr{EMG} آموزش ساخته شده است. در آخر یک \lr{decoder} خطی برای هر عضله با خروجی \lr{softplus} (که همیشه مثبت است) $K_{base}$ را می‌دهد.
\item \textbf{\lr{history gain}.} یک شبکه‌ی کوچک ($8\to16\to5$) روی ویژگی‌های تاریخچه، برای هر عضله یک ضریب $g=2^{\tanh(\cdot)}$ بین \lr{۰٫۵} و ۲ می‌سازد. در شروع آموزش $g=1$ است. خروجی نهایی $K=K_{base}\odot g$ است. این ضریب به مدل اجازه می‌دهد مثلاً بعد از چند بلندکردن سنگین، سطح فعالیت را بالاتر پیش‌بینی کند.
\end{enumerate}

\subsection{جریان \lr{EEG}}

این جریان فقط از \lr{EEG}، \lr{E} را می‌سازد:
\begin{enumerate}
\item \textbf{\lr{embedding} الکترود.} یک \lr{MLP} مشترک چهار \lr{band power} هر الکترود را به یک بردار ۳۲ بعدی تبدیل می‌کند. یک بردار آموختنی هم به هر الکترود اضافه می‌شود تا مدل بداند این کدام الکترود است. در آموزش، هر الکترود با احتمال \lr{۰٫۲} حذف می‌شود تا مدل به یک الکترود وابسته نشود.
\item \textbf{\lr{electrode graph}.} دو بلوک توجه بین ۳۲ الکترود، برای هر \lr{token} جداگانه. ماتریس ثابت $P_{ij}=-d_{ij}/\operatorname{median}(d)$ است که $d_{ij}$ فاصله‌ی دو الکترود روی سر (مختصات استاندارد ۱۰-۰۵) است. پس الکترودهای نزدیک به هم از ابتدا تمایل بیشتری به تبادل اطلاعات دارند.
\item \textbf{\lr{attention pooling}.} ۳۲ الکترود باید به یک بردار برای هر \lr{token} تبدیل شوند. یک پرسش آموختنی به همه‌ی الکترودها نگاه می‌کند و میانگین وزن‌دار آن‌ها را می‌گیرد: $p_t=\sum_e \pi_{t,e}\,n_{t,e}$. وزن $\pi_{t,e}$ می‌گوید الکترود $e$ در لحظه‌ی $t$ چقدر در این بردار سهم دارد.
\item \textbf{\lr{temporal encoder}.} ۲ بلوک توجه روی ۴۰۰ \lr{token}.
\item \textbf{سر عضله.} مانند جریان سینماتیک ولی با وزن‌های جدا؛ خروجی $E$ است.
\end{enumerate}

\subsection{\lr{fusion}}

\lr{fusion} از \lr{EEG} یک اصلاح کوچک برای \lr{K} می‌سازد. \lr{tokens} سینماتیک به‌عنوان پرسش به \lr{tokens} \lr{EEG} نگاه می‌کنند (\lr{cross-attention})؛ یعنی هر لحظه‌ی حرکت می‌پرسد «کدام لحظه‌های \lr{EEG} برای من مهم است؟». سپس یک لایه‌ی خطی برای هر عضله یک عدد اصلاحی می‌سازد:
\begin{equation}
F_{eeg}=\gamma\odot W\,\mathrm{CrossAttn}\big(\mathrm{sg}[Z^{k}],\,Z^{e}\big),\qquad F=\mathrm{relu}\big(\mathrm{sg}[K]+F_{eeg}\big).
\end{equation}
$\mathrm{sg}$ (\lr{stop-gradient}) یعنی \lr{fusion} مقدار \lr{K} و \lr{tokens} سینماتیک را می‌خواند اما آن‌ها را تغییر نمی‌دهد؛ پس \lr{K} فقط با خطای خودش آموزش می‌بیند. $\gamma\in\mathbb{R}^{5}$ یک \lr{gate} برای هر عضله است که از صفر شروع می‌شود و \lr{weight decay} ندارد. اگر \lr{EEG} چیز مفیدی به \lr{K} اضافه نکند، $\gamma$ نزدیک صفر می‌ماند و در نتیجه $F=K$. کل مدل ۳۳۷٬۷۱۲ پارامتر دارد.

\subsection{\lr{loss function}}

هر سه خروجی \lr{loss function} یکسانی در حوزه‌ی جذر دارند و \lr{loss} کل جمع آن‌هاست: $L=L_K+L_E+L_F$. برای هر خروجی $o$:
\begin{equation}
L_o=\underbrace{\mathrm{Huber}_{0.5}(\hat u_o,u)}_{\text{main},\,1.0}
+0.3\underbrace{\big(1-\mathrm{CCC}(\hat u_o,u)\big)}_{\text{ccc}}
+0.5\underbrace{\frac{1}{|\mathcal{H}|}\sum_{k\in\mathcal{H}}\Big|\overline{\hat u_o}_{k}-\overline{u}_{k}\Big|}_{\text{lift}}
+0.1\underbrace{\operatorname{mean}\big|\Delta\hat u_o-\Delta u\big|}_{\text{diff}}.
\end{equation}
\begin{itemize}
\item \textbf{\lr{main}:} خطای نقطه‌به‌نقطه با تابع هوبر.
\item \textbf{\lr{ccc}:} شکل و اندازه‌ی کلی منحنی در پنجره.
\item \textbf{\lr{lift}:} برای هر نگه‌داشتن درون پنجره ($\mathcal{H}$، با دست‌کم ۵ \lr{token}) میانگین پیش‌بینی را با میانگین واقعی مقایسه می‌کند. این جزء مستقیم همان خطایی را می‌سنجد که در مدل پایه به وزن جسم بستگی داشت.
\item \textbf{\lr{diff}:} تغییر لحظه‌به‌لحظه، تا بالا و پایین رفتن‌ها هم‌زمان باشند.
\end{itemize}
وزن اضافه‌ی قله، گران‌تر بودن کم‌برآوردی و جزء استراحت عمداً حذف شدند.

\subsection{آموزش}

هر جریان \lr{optimizer} خودش را دارد: یک \lr{AdamW} برای پارامترهای \lr{kin.*} و یکی برای \lr{eeg.*} و \lr{fusion.*} (\lr{gate} بدون \lr{weight decay}). تنظیم‌های مشترک: \lr{learning rate} $10^{-3}$، \lr{weight decay} \lr{۰٫۰۱}، ۵ دوره \lr{warm-up} و سپس \lr{cosine decay}، \lr{gradient clipping} \lr{۱٫۰} برای هر جریان جداگانه، \lr{batch size} ۳۲ و حداکثر ۱۵۰ دوره. پنجره‌های آموزش ۸ ثانیه با فاصله‌ی ۲ ثانیه از هم هستند و جای شروعشان در هر دوره به‌طور تصادفی کمی جابه‌جا می‌شود.

پیش از آموزش، چند گام یک بار با \lr{mixed precision} و یک بار بدون آن اجرا و زمان‌شان مقایسه می‌شود؛ هر کدام سریع‌تر بود انتخاب می‌شود. روی \lr{GTX 1660 Ti} \lr{mixed precision} حدود سه برابر کندتر بود و دقت کامل انتخاب شد.

\paragraph{\lr{freezing} هر جریان.} دو جریان با سرعت متفاوتی یاد می‌گیرند. اگر یکی زودتر به بهترین نقطه‌اش برسد و آموزش ادامه پیدا کند، ممکن است بدتر شود. برای جلوگیری از این، هر جریان جداگانه زیر نظر است. خطای \lr{validation} هر جریان با میانگین سه‌دوره‌ای صاف می‌شود: $L_K$ برای سینماتیک و $L_E+L_F$ برای \lr{EEG} و \lr{fusion}.

شکل~\ref{fig:freeze} قاعده را نشان می‌دهد. اگر خطای یک جریان در ۳۰ دوره‌ی پشت سر هم بهتر نشود، آن جریان به وزن‌های بهترین دوره‌اش برمی‌گردد و \lr{freeze} می‌شود: دیگر \lr{loss} برایش حساب نمی‌شود، \lr{optimizer} آن را به‌روز نمی‌کند و در حالت ارزیابی می‌ماند. جریان دیگر به آموزش ادامه می‌دهد.

یک نکته‌ی ظریف: \lr{F} روی \lr{K} ساخته می‌شود. وقتی \lr{K} \lr{freeze} می‌شود و به وزن‌های قدیمی‌ترش برمی‌گردد، \lr{F} هم کمی عوض می‌شود. برای همین شمارش صبر \lr{EEG} و \lr{fusion} از صفر شروع می‌شود تا بهترین دوره‌ی آن‌ها با همان \lr{K} نهایی سنجیده شود. آموزش وقتی تمام می‌شود که هر دو جریان \lr{freeze} شده باشند یا به ۱۵۰ دوره برسد.

\begin{figure}[htbp]
\centering
% Freeze rule at the end of every epoch.
\begin{latin}
\begin{tikzpicture}[node distance=7mm and 9mm, font=\footnotesize]
\node[box, text width=5.6cm] (ep) {end of an epoch: smoothed validation loss of each stream ($L_K$; $L_E+L_F$)};
\node[box, below=of ep, text width=5.6cm] (kq) {K not frozen and no new best for 30 epochs?};
\node[kbox, right=of kq, text width=5.2cm] (kf) {restore the best K weights and freeze K;\\ restart the patience of EEG + fusion\\ (unfreeze them if they were frozen)};
\node[box, below=10mm of kq, text width=5.6cm] (eq) {EEG + fusion not frozen and no new best for 30 epochs?};
\node[ebox, right=of eq, text width=5.2cm] (ef) {restore the best EEG + fusion weights\\ and freeze them};
\node[box, below=10mm of eq, text width=5.6cm] (both) {both streams frozen, or epoch 150 reached?};
\node[result, right=of both, text width=5.2cm] (stop) {stop: best.pt = best K + best EEG and fusion};
\node[box, below=of both, text width=5.6cm] (next) {next epoch: only the streams that are not frozen are trained};
\draw[arr] (ep)--(kq);
\draw[arr] (kq)--node[lab, above]{yes}(kf);
\draw[arr] (kq)--node[lab, right]{no}(eq);
\draw[arr] (kf.south)|-($(kq.south)!0.5!(eq.north)$);
\draw[arr] (eq)--node[lab, above]{yes}(ef);
\draw[arr] (eq)--node[lab, right]{no}(both);
\draw[arr] (ef.south)|-($(eq.south)!0.5!(both.north)$);
\draw[arr] (both)--node[lab, above]{yes}(stop);
\draw[arr] (both)--node[lab, right]{no}(next);
\draw[arr] (next.west)--++(-0.5,0)|-(ep.west);
\end{tikzpicture}
\end{latin}

\caption{قاعده‌ی \lr{freezing} در پایان هر دوره. وضعیت \lr{freezing} در \lr{last.pt} ذخیره می‌شود و اگر آموزش قطع و دوباره شروع شود، حفظ می‌شود.}
\label{fig:freeze}
\end{figure}

\subsection{پیش‌بینی روی یک سری کامل}

مدل هر بار ۸ ثانیه را می‌بیند، اما یک سری چند دقیقه طول می‌کشد. برای پیش‌بینی کل سری، پنجره‌های ۸ ثانیه‌ای با فاصله‌ی ۴ ثانیه روی سری لغزانده می‌شوند و از هر پنجره فقط ۴ ثانیه‌ی وسطش نگه داشته می‌شود (شکل~\ref{fig:stitch}). دلیلش این است که \lr{tokens} نزدیک لبه‌ی پنجره بافت کمتری دارند و پیش‌بینی‌شان کم‌دقت‌تر است. اولین و آخرین پنجره لبه‌ی بیرونی‌شان را هم نگه می‌دارند تا هیچ بخشی از سری جا نماند. در آخر، \lr{tokens} ۵۰ هرتز با درون‌یابی اسپلاین مکعبی به ۵۰۰ هرتز برمی‌گردند و با پوش واقعی ۵۰۰ هرتز مقایسه می‌شوند. همین تابع در ارزیابی، تفسیر و \lr{triptych figures} به کار می‌رود.

\begin{figure}[htbp]
\centering
% Stitched inference over one series.
\begin{latin}
\begin{tikzpicture}[x=1.45cm, y=0.55cm, font=\footnotesize]
\foreach \i/\s in {0/0,1/4,2/8,3/12}{
  \draw[semithick] (\s/4,-\i) rectangle ++(2,0.6);
  \fill[black!25] (\s/4+0.5,-\i) rectangle ++(1,0.6);
  \node[lab, anchor=east] at (-0.1,-\i+0.3) {window \the\numexpr\i+1\relax};
}
\fill[black!25] (0,0) rectangle (0.5,0.6);
\fill[black!25] (3.5,-3) rectangle (5,-2.4);
\draw[arr] (-0.2,-4.1)--(5.4,-4.1);
\foreach \x/\t in {0/0,1/4,2/8,3/12,4/16,5/20}{\draw (\x,-4.0)--(\x,-4.2) node[lab, below]{\t};}
\node[lab] at (2.5,-5.4) {time in the series (s)};
\node[lab, anchor=west, align=left] at (5.5,-1.2) {grey: the part kept from each\\ 8 s window (its central 4 s;\\ the first and last windows\\ also keep their outer edge)};
\end{tikzpicture}
\end{latin}

\caption{پیش‌بینی پیوسته روی یک سری: پنجره‌های ۸ ثانیه‌ای با فاصله‌ی ۴ ثانیه، و نگه‌داشتن ۴ ثانیه‌ی وسط هر پنجره.}
\label{fig:stitch}
\end{figure}
